\section{Problem formulation}
\label{sec:problem}

We consider an autonomous vehicle navigating a roadway over planning horizon $T_{\mathrm{pred}}$ in the presence of $k(t)$ static and dynamic obstacles. Given vehicle kinematic states and obstacle motion predictions, our objective is to systematically construct the spatial and temporal constraint interface $(\mathcal{B}_l,\mathcal{B}_s)$ for two-stage path--velocity optimization without altering the structure of downstream convex quadratic programs. Following standard formulations in autonomous vehicle motion planning, we focus on behavioral- and trajectory-level optimization; low-level tracking controllers are assumed to track reference trajectories, while communication delays and low-level powertrain dynamics are decoupled from the upper-level trajectory generation.

\subsection{Path--velocity decomposition interface}

We formulate vehicle motion planning in a road-aligned Frenet coordinate frame defined with respect to a designated smooth reference line. In this frame, the longitudinal coordinate $s$ denotes the arc length along the reference line, representing the vehicle's forward travel distance, while the lateral coordinate $l$ denotes the signed orthogonal displacement from the reference line.

Under the two-stage path--velocity decomposition framework, the complete spatio-temporal motion planning problem is decoupled into two sequential stages:
first, solving for a spatial path profile $l(s)$ that determines lateral maneuvering and obstacle avoidance along the longitudinal station;
second, solving for a longitudinal trajectory profile $s(t)$ that governs the vehicle's forward position, speed, and acceleration over the planning horizon $t \in [0, T_{\mathrm{pred}}]$.

To render these continuous profiles tractable for real-time numerical optimization, the curves are discretized over spatial stations $s_j$ for $j=0,\dots,K$ and temporal steps $t_r$ for $r=0,\dots,N$. Let the discrete path state vector be $\mathbf{x}_p=(l_j,l'_j,l''_j)_{j=0}^{K}$, where $l'_j = \mathrm{d}l/\mathrm{d}s$ and $l''_j = \mathrm{d}^2l/\mathrm{d}s^2$ represent the lateral heading deflection and curvature rate, respectively. Let the discrete longitudinal trajectory state vector be $\mathbf{x}_v=(s_r,v_r,a_r)_{r=0}^{N}$, where $v_r = \dot{s}(t_r)$ and $a_r = \ddot{s}(t_r)$ denote the longitudinal velocity and acceleration. The two planning stages are sequentially solved via convex quadratic programs:
\begin{align}
 \mathbf{x}_p^*={}&\arg\min_{\mathbf{x}_p}
 \frac{1}{2}\mathbf{x}_p^{\mathsf T}H_p\mathbf{x}_p+q_p^{\mathsf T}\mathbf{x}_p,\\[-2pt]
 &\quad\text{s.t. }A_p\mathbf{x}_p=b_p,\quad l_j^-\leq l_j\leq l_j^+, \label{eq:path_qp_en}\\
 \mathbf{x}_v^*={}&\arg\min_{\mathbf{x}_v}
 \frac{1}{2}\mathbf{x}_v^{\mathsf T}H_v\mathbf{x}_v+q_v^{\mathsf T}\mathbf{x}_v,\\[-2pt]
 &\quad\text{s.t. }A_v\mathbf{x}_v=b_v,\quad s_r^{lb}\leq s_r\leq s_r^{ub},\notag\\[-2pt]
 &\qquad 0\leq v_r\leq v_{\max},\quad a_{\min}\leq a_r\leq a_{\max}. \label{eq:speed_qp_en}
\end{align}
Here, $H_p$ and $H_v$ encode trajectory smoothness costs, while $A_p$ and $A_v$ enforce piecewise-jerk kinematic transition equalities between adjacent discretization points. The operational environment is exposed to the solvers strictly through the spatial path boundaries $\mathcal{B}_l=\{[l_j^-,l_j^+]\}_{j=0}^K$ and temporal position boundaries $\mathcal{B}_s=\{[s_r^{lb},s_r^{ub}]\}_{r=0}^N$. MIKU reconstructs these constraint interfaces while preserving the objective matrices, kinematic equalities, and convex QP solver formulations.

\subsection{Center-coordinate obstacle model}

For obstacle $i$, let $O_i(t)=[s_i^-(t),s_i^+(t)]\times[l_i^-(t),l_i^+(t)]$ denote its physical bounding box. With road limits $[l_{\mathrm{road}}^-,l_{\mathrm{road}}^+]$, vehicle width $W_{\mathrm{ego}}$, and roadside margin $d_{\mathrm{road}}$, the feasible interval for the vehicle center is
\begin{equation}
 [L^-,L^+]=[l_{\mathrm{road}}^-+W_{\mathrm{ego}}/2+d_{\mathrm{road}},\ l_{\mathrm{road}}^+-W_{\mathrm{ego}}/2-d_{\mathrm{road}}].
 \label{eq:center_road_en}
\end{equation}
The obstacle occupancy interval is expanded once by half the vehicle width and a threat-dependent buffer $d_{\mathrm{buf},i}$:
\begin{equation}
 [u_i,v_i]=[l_i^- -W_{\mathrm{ego}}/2-d_{\mathrm{buf},i},\ l_i^+ +W_{\mathrm{ego}}/2+d_{\mathrm{buf},i}].
 \label{eq:center_forbidden_en}
\end{equation}
All quantities in Equations~\ref{eq:center_road_en}--\ref{eq:center_forbidden_en} are defined with respect to the vehicle center, ensuring that geometric dimensions are not subtracted repeatedly in downstream modules.

Assigning each of the $k$ active obstacles at a given longitudinal station to the left or right side via decision vector $\mathbf{d}$ yields
\begin{equation}
 l^-(\mathbf{d})=\max\left(L^-,\max_{i\in\mathcal{L}}v_i\right),\qquad
 l^+(\mathbf{d})=\min\left(L^+,\min_{i\in\mathcal{R}}u_i\right),
 \label{eq:center_band_en}
\end{equation}
with empty index sets omitted. The resulting lateral band is feasible when the available clearance satisfies $W(\mathbf{d})=l^+(\mathbf{d})-l^-(\mathbf{d})\geq\varepsilon$ for a minimum width threshold $\varepsilon$.

\subsection{Baseline failure mechanisms and evaluation criteria}

A time-blind path stage uses the full prediction horizon union
\begin{equation}
 \overline{O}_i=\bigcup_{t\in[0,T_{\mathrm{pred}}]}O_i(t),
 \label{eq:prob_union_en}
\end{equation}
treating spatial occupancies across different time instances as simultaneously blocked. Furthermore, greedy obstacle-wise side decisions can push the upper and lower bounds in Equation~\ref{eq:center_band_en} past each other, resulting in $l^+(s)<l^-(s)$ and falsely declaring an artificial blockage even when a collision-free spatio-temporal passage exists.

An episode is deemed successful when the optimized trajectory composed of $\mathbf{x}_p^*$ and $\mathbf{x}_v^*$ remains collision-free throughout $[0,T_{\mathrm{pred}}]$ and reaches $s\geq s_{\mathrm{end}}-1$~m, with the complete operational evaluation metrics detailed in Section~\ref{sec:experimental_design}.
